% ==============================================================================
% CAPÍTULO 5 - CONSIDERAÇÕES FINAIS
% ==============================================================================

\section{Considerações Finais}
\label{sec:consideracoes}

Este trabalho perguntou se um \textit{harness} em nível de repositório compensa
a escala do modelo em agentes de código. Para respondê-la, construiu um
\textit{benchmark} de ablação com 14 tarefas derivadas de \textit{pull requests}
reais de uma aplicação Laravel privada, avaliadas por testes ocultos; um
orquestrador que executou seis modelos em quatro níveis cumulativos de
\textit{harness}, em \textit{sandbox} e com registro reprodutível; e uma análise
de qualidade de código em três camadas sobre 309 execuções e 836 revisões de
juízes. Os cinco objetivos específicos foram atingidos: o \textit{benchmark}, os
níveis e a execução foram realizados como planejado, com a redução de três
repetições para uma registrada e justificada; a análise de qualidade e a
análise estatística pareada foram conduzidas e seus resultados discutidos.

A resposta à pergunta é dupla e depende da métrica. Na taxa de resolução, o
\textit{harness} não compensou a escala: quatro dos seis modelos já resolviam
95\% ou mais das tarefas sem nenhuma estrutura de apoio, nenhum nível alterou o
\textit{pass@1} de nenhum modelo (McNemar exato, $p = 1{,}0$ em 70 pares por
nível) e cada execução ficou mais cara em turnos e \textit{tokens}. Entre
modelos, o \textit{pass@1} variou de 11\% a 100\%; entre níveis de um mesmo
modelo, no máximo 7 pontos percentuais. Na qualidade do código, o
\textit{harness} teve efeito consistente: a nota geral dos juízes subiu 0,31 a
0,34 ponto, a nota de convenções 0,46 a 0,52 e a de testes 0,60 a 0,87, com
cerca de dez casos de teste a mais por execução, menos erros novos de análise
estática e menor complexidade, e sem alteração em requisitos, segurança ou
escopo. O efeito concentrou-se nos modelos de capacidade intermediária
(DeepSeek v4.1 flash e Muse Spark 1.3), foi pequeno nos modelos de ponta, que
já partem alto, e nulo nos modelos mais fracos (gpt-oss-20b e Beam), que não
seguem as regras. A primeira camada, as regras sempre ativas, respondeu pela
maior parte do efeito e do custo; as \textit{skills} acrescentaram pouco e a
configuração de LSP nada mensurável. Em síntese, nesta amostra a escala
explica se a tarefa é resolvida e o \textit{harness} explica parte de como o
código é entregue.

As contribuições são o \textit{benchmark} e seus enunciados, o orquestrador, o
pacote de análise de qualidade com a rubrica e as regras de convenção, e os
dados de todas as execuções e revisões, disponíveis publicamente, com exceção
dos \textit{commits} do repositório-sujeito, que é privado. Do ponto de vista
metodológico, o trabalho mostra que um oráculo de testes, sozinho, não
distingue entregas que um revisor distinguiria, e que juízes LLM de três
famílias, com evidência obrigatória, controle de autopreferência e validação
cruzada por medidas objetivas, podem preencher essa lacuna com concordância
substancial.

As limitações são as de um experimento de uma repetição, 14 tarefas de um
único repositório privado e efeito de teto: o desenho só detectaria efeitos
grandes sobre o \textit{pass@1}, e nada aqui se generaliza para outros
repositórios, linguagens ou \textit{harnesses}. O oráculo foi parcialmente
escrito com IA, os juízes são LLMs sem calibração humana, as regras de
convenção são heurísticas, o Beam ficou incompleto e as comparações múltiplas
não foram corrigidas.

Como trabalhos futuros, propõem-se: executar mais repetições por célula, para
estimar a variância de um mesmo modelo e ganhar poder estatístico; selecionar
tarefas mais difíceis, fora do teto, nas quais um ganho de \textit{pass@1} seja
mensurável; repetir o protocolo em outros repositórios e pilhas, com seus
próprios \textit{harnesses}, para testar a generalidade dos achados; fazer a
ablação de partes do \texttt{AGENTS.md} (apenas a regra de testes, apenas as
convenções, apenas a política de segurança), para identificar quais frases
carregam o efeito; e calibrar os juízes LLM com revisores humanos em uma
amostra das mudanças, para medir a validade das notas diretamente. Uma rodada
própria para os modelos locais, com o Qwen retirado desta, e a conclusão da
matriz do Beam completariam o quadro dos modelos de menor capacidade, onde a
hipótese \textit{scaffolding over scale} tem mais a provar.
